% Recovered verbatim from chapters_6_10/chapters_6_10.tex
% Canonical chapter 6
\chapter{ER Network Formalism}

\section{Planck-Scale Adjacency Graph}

We model microscopic spatial adjacency not by metric distance but by 
an Einstein--Rosen (ER) connectivity graph:
\[
W_{ij} \ge 0,\qquad W_{ij}=W_{ji},
\]
where $i,j$ label Planck-scale sites.  
Edges represent wormhole-like entanglement connections.

\section{ER Contribution to Geometry}

We define an effective ER metric component:
\[
g^{(ER)}_{\mu\nu} = \lambda 
\sum_{i,j} W_{ij}\, \xi^{(i)}_\mu \xi^{(j)}_\nu,
\]
where $\xi^{(i)}_\mu$ projects fields onto site neighborhoods.

\section{Entanglement and Curvature}

The Bianchi-compatibility condition suggests:
\[
R_{\mu\nu} \sim \nabla_\mu\nabla_\nu S_{\mathrm{ER}},
\]
where $S_{\mathrm{ER}}$ is ER-network entanglement entropy.


